\section{What It Costs, and What Is Still Missing}
\label{sec:ble-costs}

A chapter about a radio stack that only reports successes would be useless. This one ends with the
bill.

\subsection{The Numbers}

\begin{table}[htbp]
    \centering
    \footnotesize
    \begin{tabular}{@{}p{0.30\textwidth}p{0.22\textwidth}p{0.40\textwidth}@{}}
        \toprule
        \textbf{Example} & \textbf{Image size} & \textbf{What it needs} \\
        \midrule
        \codeinline{ble\_beacon} & ~5.9 KB & three blobs, the stubs, a delay \\
        \codeinline{rail\_version} & ~1.3 KB & one blob, the stubs \\
        \codeinline{ble\_gatt} & does not link & the above plus PSA crypto \\
        \bottomrule
    \end{tabular}
    \caption{Measured with elf-load-range.py. The application is small; the blobs are
    not, and they are what fills flash.}
    \label{tab:ble-sizes}
\end{table}

The three Bluetooth blobs are on the order of a megabyte together. On a part with 1.5 MB of flash and
a 24 KiB bootloader, that is affordable once — and it is the reason a project chooses BLE for a
product rather than adding it ``just in case''.

\subsection{The Gap That Must Be Named: PSA Crypto}

A connectable stack references the PSA crypto API, and this SDK build does not ship it in the blobs.
The eight symbols are:

\begin{minted}{text}
psa_crypto_init      psa_destroy_key       psa_import_key    psa_mac_abort
psa_mac_sign_finish  psa_mac_sign_setup    psa_mac_update    psa_reset_key_attributes
\end{minted}

They come from Silicon Labs' \codeinline{libpsa\_crypto} and \codeinline{libmbedtls} builds, which are
separate artefacts with their own licence terms. Until they are vendored, a GATT example is a staged
file rather than a working one — and this book says so instead of shipping it as an example that
``should work''.

A project that does not need bonding can sometimes avoid the crypto path, but not by accident: the
link tells you whether it has.

\subsection{The Three Things to Verify Before Believing a BLE Build}

\begin{enumerate}
    \item \textbf{Version match.} The header and the blobs are from the same SDK. Check a symbol the
          blob exports (\codeinline{arm-none-eabi-nm libbluetooth_host*.a | grep sl\_bt\_advertiser})
          against the generated bindings.
    \item \textbf{Link, not compile.} A \codeinline{-sys} crate that compiles proves nothing about
          the blobs. Build one example that links.
    \item \textbf{One real advertisement.} A phone that sees the name is the only proof that the
          antenna, the crystal and the stack configuration agree. Everything below that is a build
          check.
\end{enumerate}

\subsection{Where This Book Stands}

\begin{minted}{text}
bindings compile                    yes, with three documented workarounds
beacon links                        yes
beacon flashed and seen by a phone  not yet -- the next measurement
GATT links                          no, PSA crypto missing
RAIL links for version only         yes; RAIL_Init needs emlib sources
\end{minted}

Stating that plainly is the point. A chapter that claimed the last four were done would be worse than
no chapter, because the next person would trust it.
